\section{Related Work}
\label{sec:sota}

The line of work that frames our contribution starts from contrastive vision--language pretraining and traces how each successive refinement---adaptation to video, specialisation to first-person footage, accommodation of graded relevance, and inference-time fusion---reshaped what is achievable on benchmarks like EK-100 MIR. CLIP~\cite{radford2021clip} established the paradigm by showing that contrastive alignment of $400$M image--text pairs produces general-purpose visual representations whose zero-shot transfer matches fully supervised baselines, displacing the need to train task-specific encoders from scratch. Bringing the paradigm to video required rethinking how individual frames combine into a clip-level signal: CLIP4Clip~\cite{luo2022clip4clip} compared three temporal aggregation strategies and found that the quality of the pretrained representation outweighs the choice of aggregation, an observation that X-CLIP~\cite{ma2022xclip} pushed further with multi-granularity contrastive learning that aligns video--sentence pairs at the global level and frame--word pairs at the local level via a cross-frame communication transformer. ViFi-CLIP~\cite{rasheed2023vificlip} closed this trajectory by showing that end-to-end fine-tuning with minimal architectural modifications already achieves competitive zero-shot generalisation, evidence that shared embedding spaces transfer well to new domains and motivating our choice of a CLIP-style dual encoder fine-tuned directly with retrieval-aware supervision in Approach~3.

These third-person results, however, do not transfer cleanly to the first-person setting that drives EK-100. Egocentric footage exhibits dynamic camera motion, high intra-class variability and a distribution very different from third-person video, and only at the scale of Ego4D~\cite{grauman2022ego4d}---$3{,}600$\,h of first-person recordings from 931 participants---did large-scale pretraining for this domain become feasible. EPIC-KITCHENS-100~\cite{damen2022rescaling} provided the cooking-domain reference benchmark and, crucially, formalised the multi-instance retrieval task with \emph{soft-label} relevance matrices, evaluated with order-sensitive metrics (nDCG~\cite{jarvelin2002ndcg}, mAP) rather than the binary R@K used in third-person retrieval; the construction of these matrices was later given a principled grounding by Wray et al.~\cite{wray2021semantic} via word2vec similarity over verb and noun classes. This data infrastructure enabled a generation of egocentric encoders that progressively closed the gap with third-person performance: EgoVLP~\cite{lin2022egovlp} introduced EgoNCE, an InfoNCE variant that weights additional positives through verb/noun matching---improving over general-domain methods on EK-100 but still relying on lexical overlap rather than on the official graded matrices; LaViLa~\cite{zhao2023lavila} fine-tuned a language model on Ego4D to generate diverse training narrations, becoming the reference video encoder for subsequent work; EgoVLPv2~\cite{pramanick2023egovlpv2} integrated cross-modal fusion directly into the backbone via interleaved cross-attention; and AVION~\cite{zhao2023avion}, building on LaViLa's encoder design, trained a ViT-L/14 backbone on a single machine in a single day, won the CVPR~2023 EK-100 challenge and, in combination with FlashAttention~\cite{dao2022flashattention} for memory-efficient computation, became the de-facto backbone for downstream retrieval work---and, in our pipeline, the visual encoder of Approach~3.

Strong encoders alone, however, do not exploit the soft-label structure of the EK-100 MIR task: the loss function decides whether the graded relevance $R$ enters the gradient signal at all. The official challenge baseline reflects this gap, combining JPoSE~\cite{wray2019jpose}---which learns separate verb, noun and action embedding spaces---with MI-MM, a multi-instance MaxMargin adaptation of MIL-NCE~\cite{miech2020milnce}; both operate with binary relevance and ignore the soft-label matrices. JPoSE itself was built on top of the simpler Multi-Modal Embedding Network (MMEN), an MLP-based projection used as a baseline in the same work and which we revisit as the third member of our score-level ensemble. The decisive step beyond binary relevance came with the Symmetric Multi-Similarity (SMS) loss of Wang, Wang and Chau~\cite{wang2024symmetricmultisimilaritylossepickitchens100}, which extends the multi-similarity formulation of~\cite{wang2019multisimilarity} so that, instead of binary positive/negative pair weights, the contribution of each pair is modulated by its soft relevance score, applied symmetrically over both modalities; built on LaViLa as the base encoder, SMS is the most explicit attempt to date to directly optimise the official challenge metric and is the loss we adopt in Approach~3. An adjacent line of research within the broader retrieval community has developed losses that approximate ranking metrics directly---SoDeep~\cite{engilberge2019sodeep} learns a sorting surrogate for differentiable mAP and Recall, and Smooth-AP~\cite{brown2020smoothap} provides a sigmoidal AP approximation---but neither family has been systematically explored in video--text retrieval with graded relevance, leaving a clear gap that our discussion in Section~\ref{sec:conclusions:future} returns to. Beyond loss design, scaling has continued to push the limits of multimodal understanding (InternVideo2~\cite{wang2024internvideo2} reaches state of the art across more than sixty benchmarks via masked video pretraining, contrastive alignment and instruction tuning), and the current EK-100 MIR public leader is CR-CLIP~\cite{he2025crclip}, which builds on the AVION dual-encoder architecture and adds a context-refinement module to reach $82.10$\,nDCG and $66.8$\,mAP.

Because each of the strands above reshapes only a single component of the retrieval pipeline---the encoder, the loss or the data---there remains a complementary line of work that operates at inference time, and which our Approach~2 explicitly tests. Score-level ensemble fusion combines similarity matrices of complementary models, typically by row-wise normalisation followed by weighted averaging, to reduce variance with no additional training cost. Re-ranking refines initial scores by exploiting the local geometry of the retrieval space: Iscen et al.~\cite{iscen2017diffusion} introduced an efficient diffusion procedure on region manifolds in which scores are propagated through a sparse k-nearest-neighbour graph, and k-reciprocal encoding~\cite{zhong2017kreciprocal} formalises mutual nearest-neighbour relations as a re-ranking signal; we adapt these ideas to the video--text setting by propagating ensemble scores through both visual and textual k-NN graphs (Eq.~\eqref{eq:rerank}). At the parameter level, \emph{model soups}~\cite{wortsman2022modelsoups} average the weights of several models fine-tuned with different hyperparameters and consistently match or exceed the best individual model with inference cost equal to that of a single one. Wang et al.~\cite{wang2024symmetricmultisimilaritylossepickitchens100} apply a closely related strategy in their winning EK-100 MIR submission---a six-model score-level ensemble that mixes Adaptive MI-MM and SMS variants trained at the same length but with different learning-rate schedules and relaxation thresholds---which we discuss explicitly in Section~\ref{sec:conclusions:future} as the natural continuation of our pipeline.

Read together, these four strands expose a structural mismatch with the EK-100 MIR setting that no single one resolves: methods designed for third-person video assume binary relevance; egocentric-specialised encoders such as EgoVLP and LaViLa only partially leverage the soft-label matrices (EgoNCE replaces them with lexical overlap, LaViLa applies standard contrastive learning); approaches that consume the soft labels directly, such as SMS, do so as a fine-tuning stage on top of binary-pretrained representations; and ensemble or re-ranking strategies have rarely been combined explicitly with the soft-label EK-100 setting. Our work bridges these gaps along a single, controlled iteration trajectory: Approach~1 reproduces the binary-relevance baseline, Approach~2 quantifies how far inference-time fusion of frozen soft-label-agnostic models can carry it, and Approach~3 fine-tunes a strong egocentric encoder with relevance-aware supervision on the same evaluation protocol, isolating the contribution of every design choice along the way.
